% Work-MAT | 2. rocnik | pripravovana hodina 39
% Návaznost: uživatelské připravy.zip (hodiny 01–20), autorská učebnice na GitHubu.
% Realisticky.cz: didaktická inspirace v posloupnosti a typech úloh; zadání jsou vlastní.
\documentclass[11pt,a4paper]{article}
\input{../../sablony/pripravy/priprava.sty}
\nastavPripravu{2. ročník}{Logaritmické rovnice}{39}{Rovnice s více logaritmy}
\begin{document}
\blok{Motivace a návaznost}
Žáci nejprve určují průnik podmínek několika logaritmů. Po sloučení získají často kvadratickou rovnici, jejíž kořeny je nutné filtrovat.
\poznamka{Orientační nabídka pro 45--55 minut: krátké zopakování a motivace, nový výklad, tři společné vzorové úlohy, samostatná práce. Vyber podle tempa třídy; zbývající zadání slouží jako rezerva.}
\blok{Výklad}
Každý argument logaritmu musí být kladný současně; následující algebraické kroky nelze použít k jejich opomenutí.
\blok{Společné příklady}
\begin{prikladReseniBox}[title={Příklad 1}]
Vyřeš, resp. urči: $\log_2 x+\log_2(x-2)=3$.
\tcblower
\[\begin{aligned}x>2 \\ \log_2(x(x-2))=3 \\ x^2-2x=8\;\Longrightarrow\;x=4\end{aligned}\]
\end{prikladReseniBox}
\begin{prikladReseniBox}[title={Příklad 2}]
Vyřeš, resp. urči: $\log_3(x-1)+\log_3(x+1)=1$.
\tcblower
\[\begin{aligned}x>1 \\ x^2-1=3 \\ x=2\end{aligned}\]
\end{prikladReseniBox}
\poznamka{Povinný krok je zápis průniku definičních oborů. Záporný algebraický kořen nevyhovuje ani když formálně řeší pomocnou rovnici.}
\blok{Samostatné procvičování}
\begin{zadaniReseniBox}[title={Samostatná práce 1}]
\begin{enumerate}[label=\alph*)]
\item $\log_{2}x+\log_{2}\left(x-1\right)=\log_{2}\left(6\right)$.
\item $\log_{3}x+\log_{3}\left(x-2\right)=\log_{3}\left(99\right)$.
\item $\log_{5}x+\log_{5}\left(x-3\right)=\log_{5}\left(16000\right)$.
\item $\log_{2}x+\log_{2}\left(x-1\right)=\log_{2}\left(6\right)$.
\item $\log_{3}x+\log_{3}\left(x-2\right)=\log_{3}\left(99\right)$.
\item $\log_{5}x+\log_{5}\left(x-3\right)=\log_{5}\left(16000\right)$.
\end{enumerate}
\tcblower
\begin{enumerate}[label=\alph*)]
\item $x=3 (x>1)$.
\item $x=11 (x>2)$.
\item $x=128 (x>3)$.
\item $x=3 (x>1)$.
\item $x=11 (x>2)$.
\item $x=128 (x>3)$.
\end{enumerate}
\end{zadaniReseniBox}
\blok{Dobrovolné rozšíření}
\begin{zadaniReseniBox}[title={Rozšíření 1}]
\begin{enumerate}[label=\alph*)]
\item $Vytvoř vlastní úlohu tohoto typu a se spolužákem porovnej podmínky a výsledek.$.
\end{enumerate}
\tcblower
\begin{enumerate}[label=\alph*)]
\item Zkontroluj definice a argumenty..
\end{enumerate}
\end{zadaniReseniBox}
\blok{Úkol na doma}
U každého logaritmu vypiš podmínky argumentu a proveď zkoušku. U rovnic vyznač případně nepřípustné kořeny, u nerovnic odpověz intervalem.
\begin{zadaniReseniBox}[title={Domácí úkol}]
\begin{enumerate}[label=\alph*)]
\item $\log_{2}x+\log_{2}\left(x-1\right)=\log_{2}\left(6\right)$.
\item $\log_{3}x+\log_{3}\left(x-2\right)=\log_{3}\left(99\right)$.
\item $\log_{5}x+\log_{5}\left(x-3\right)=\log_{5}\left(16000\right)$.
\item $\log_{2}x+\log_{2}\left(x-1\right)=\log_{2}\left(6\right)$.
\item $\log_{3}x+\log_{3}\left(x-2\right)=\log_{3}\left(99\right)$.
\item $\log_{5}x+\log_{5}\left(x-3\right)=\log_{5}\left(16000\right)$.
\item $\log_{2}x+\log_{2}\left(x-1\right)=\log_{2}\left(6\right)$.
\item $\log_{3}x+\log_{3}\left(x-2\right)=\log_{3}\left(99\right)$.
\end{enumerate}
\tcblower
\begin{enumerate}[label=\alph*)]
\item $x=3 (x>1)$.
\item $x=11 (x>2)$.
\item $x=128 (x>3)$.
\item $x=3 (x>1)$.
\item $x=11 (x>2)$.
\item $x=128 (x>3)$.
\item $x=3 (x>1)$.
\item $x=11 (x>2)$.
\end{enumerate}
\end{zadaniReseniBox}
\poznamka{Domácí úkol obsahuje 8 položek; předpokládaná domácí práce přibližně 15--30 minut. Po větších blocích jsou vědomě promíchané dříve probrané typy.}
\end{document}
